problem:  
Let \( (M^{2n}, g) \) be a compact, simply connected Riemannian manifold with **nonnegative sectional curvature**, admitting an effective, isometric torus action of rank \(k\) that is **slice‑maximal** in the precise sense of Galaz‑García–Kerin–Searle (GGKS); that is, equality holds in the GGKS rank–dimension inequality.  
By their theorem, \(M\) is equivariantly **rationally homotopy equivalent** to a torus orbifold \(\mathcal{O}_P\) determined by labeled simple polytope data \(P\).

Consider the Borel fibration  
\[
M \hookrightarrow M_{T^k} \to BT^k
\]
whose Leray–Serre spectral sequence over \(\mathbb Z\) has \(E_2^{p,q} = H^p(BT^k; H^q(M;\mathbb Z))\).  
In general, this spectral sequence may not collapse at \(E_2\) over \(\mathbb Z\), even when it collapses rationally, because torsion in \(H^*(M;\mathbb Z)\) may drive non‑zero differentials.  

**Open problem:**  
Under the above hypotheses, does nonnegative curvature force **integral collapse** of this spectral sequence, i.e.  
\[
E_2 = E_\infty, \quad \text{so that } H^*_{T^k}(M;\mathbb Z) \text{ is a free } H^*(BT^k;\mathbb Z)\text{‑module}?
\]  
Equivalently: can nonnegative sectional curvature eliminate all torsion differentials in the Borel spectral sequence for slice‑maximal torus actions?  
If not, can one construct a counterexample—namely, a nonnegatively curved, slice‑maximal manifold whose Borel spectral sequence exhibits non‑trivial integral differentials driven by torsion?

Why is it a "good" problem:  
This problem captures, in a compact and verifiable form, a sharp interface between **geometry, topology, and algebra**. The question is simple to state in terms of spectral sequence collapse yet probes principal structural constraints linking curvature to torsion in equivariant cohomology.  

- **Bridges multiple frameworks:** It unites Riemannian comparison geometry (nonnegative curvature), transformation group theory (slice‑maximal torus actions), and equivariant algebraic topology (integral formality and torsion phenomena).  
- **Conceptual depth:** While rational equivariant formality is known (by GGKS), integral formality under curvature constraints is entirely unexplored. The question asks whether geometric positivity induces algebraic purity—in other words, whether curvature “kills torsion.”  
- **Test of rigidity vs. flexibility:** A positive answer would reveal a new rigidity principle for curved, symmetric spaces and forge deeper ties between curvature and the arithmetic of cohomology rings. A counterexample would expose unexpected torsion mechanisms compatible with nonnegative curvature.  

Because it isolates a precise algebraic condition (spectral sequence degeneracy) under a geometric hypothesis (nonnegative curvature) and connects known rational models to new integral phenomena, this problem embodies the essential qualities of a good research‑level problem: clarity, simplicity of statement, and profound potential consequences across fields.